\section{Strategische Ma{\ss}nahmen im Gesch\"aftsbericht 2024}
\label{sec:strategie}

Dieser Abschnitt gibt die im Gesch\"aftsbericht 2024 berichteten Ma{\ss}nahmen
wieder. Die Darstellung ist rein sachlich; die Bewertung dieser Ma{\ss}nahmen
erfolgt in Abschnitt~\ref{sec:optionen}.

Aumann gliedert das Gesch\"aft in zwei Segmente: E-Mobility und Next
Automation, vormals Classic.\quelleGB{2024}{9}\quelleMerken{p10-gb2024-9} Das Segment E-Mobility machte
im Gesch\"aftsjahr 2024 rund \Pct{\AnteilEMobility} des Gesch\"afts
aus.\quelleGB{2024}{10}\quelleMerken{p10-gb2024-10}

\subsection{Wachstumssegment E-Mobility}
\label{subsec:e-mobility}

Aumann bezeichnet sich als global t\"atiger Hersteller von Spezialmaschinen
und hochautomatisierten Produktionslinien mit strategischem Schwerpunkt auf
der Elektromobilit\"at.\quelleErneut{p10-gb2024-9} Das Leistungsspektrum reicht von
modularen Produktionszellen \"uber Prozessl\"osungen (etwa f\"ur Wickel- und
Beschichtungstechnik) bis zu schl\"usselfertigen
Gro{\ss}serienanlagen.\quelleErneut{p10-gb2024-9}

Die Fertigungsl\"osungen adressieren E-Traktionsantriebe, Batteriesysteme,
Leistungselektronik sowie Elektronik- und Sensorkomponenten und
Fahrwerksteile.\quelleErneut{p10-gb2024-10} Kunden setzen die Anlagen f\"ur die
Serienfertigung vollelektrischer und hybrider Fahrzeugantriebe sowie zur
Herstellung von Energiespeichersystemen, E-Traktionsmotoren und elektrischen
Nebenaggregaten ein.\quelleErneut{p10-gb2024-10}

Im Gesch\"aftsjahr 2024 stieg der Segmentumsatz um
\Pct{\UmsatzEMobilityDeltaZFIV} auf \MioEUR{\UmsatzEMobilityZFIV}; das
Segment-EBITDA erh\"ohte sich von \MioEUR{\EbitdaEMobilityZFIII} auf
\MioEUR{\EbitdaEMobilityZFIV}, was einer Marge von
\Pct{\EbitdaMargeEMobilityZFIV} entspricht.\quelleGB{2024}{13}\quelleMerken{p10-gb2024-13}

\subsection{Neuausrichtung des Segments Next Automation}
\label{subsec:next-automation}

Mit der Ver\"offentlichung der Quartalsmitteilung zum 30.~September 2024 am
14.~November 2024 kommunizierte Aumann die Neuausrichtung des vormaligen
Segments Classic.\quelleErneut{p10-gb2024-10} Dieses fasste bis dahin Produktionslinien
f\"ur Komponenten konventionell angetriebener Fahrzeuge und f\"ur
Anwendungen au{\ss}erhalb der Automobilindustrie zusammen. Au{\ss}erhalb der
Automobilindustrie waren bereits Referenzprojekte umgesetzt worden, unter
anderem automatisierte Fertigungslinien f\"ur Haushaltsger\"ate, industrielle
Elektromotoren, Photovoltaikmodule, Wasserstoff-Elektrolyseure und
Windkraftanlagen.\quelleErneut{p10-gb2024-10}

Das Segment wurde in \enquote{Next Automation} umbenannt und wird k\"unftig
auf Automatisierungsl\"osungen f\"ur Anwendungen wie Clean Tech, Aerospace und
Life Sciences ausgerichtet.\quelleErneut{p10-gb2024-10} Der Bericht nennt als Vorbild
die Entwicklung des Segments E-Mobility.\quelleErneut{p10-gb2024-10}

\begin{table}[htbp]
\centering
\caption{Segment Next Automation: Umsatz und EBITDA in Mio.\,\euro,
Marge in Prozent.}
\label{tab:next-automation}
\begin{tabular}{@{}lrrr@{}}
\toprule
\textbf{Gr\"o{\ss}e} & \textbf{2023} & \textbf{2024} & \textbf{2025}\\
\midrule
Umsatz       & \num{\UmsatzNextAutoZFIII} & \num{\UmsatzNextAutoZFIV} & \num{\UmsatzNextAutoZFV}\\
EBITDA       & \num{\EbitdaNextAutoZFIII} & \num{\EbitdaNextAutoZFIV} & \num{\EbitdaNextAutoZFV}\\
EBITDA-Marge & \Pct{\EbitdaMargeNextAutoZFIII} & \Pct{\EbitdaMargeNextAutoZFIV}
             & \Pct{\EbitdaMargeNextAutoZFV}\\
\bottomrule
\end{tabular}
\end{table}

\noindent Der Auftragseingang des Segments belief sich 2024 auf
\MioEUR{\AuftragseingangNextAutoZFIV}\quelleErneut{p10-gb2024-13} und 2025 auf
\MioEUR{\AuftragseingangNextAutoZFV}.\quelleGB{2025}{13}

\subsection{Batterie- und Brennstoffzellen-Produktionstechnik}
\label{subsec:batterie}

Der Bericht ordnet dem Feld der Elektromobilit\"at ausdr\"ucklich Energie\-speicher-
und Energiewandlungssysteme aus Batterie und Brennstoffzelle
zu.\quelleGB{2024}{9}\quelleMerken{p11-gb2024-9} Zum Leistungsangebot z\"ahlen Laminier- und
Beschichtungsanlagen f\"ur die Elektroden- und MEA-Fertigung
(Membran-Elektroden-Einheit).\quelleGB{2024}{10}\quelleMerken{p11-gb2024-10}

Die Produktionstechnik f\"ur die Brennstoffzellen- und Elektrolyseurfertigung
war ein zentraler Schwerpunkt von Forschung und
Entwicklung.\quelleGB{2024}{13}\quelleMerken{p11-gb2024-13} Weiterentwickelt wurden zudem Verfahren zur
Herstellung von Batterie-Elektrodenfolien.\quelleGB{2024}{14}

Als j\"ungste Akquisition nennt der Bericht die Aumann Lauchheim GmbH, deren
Post-Merger-Integration und Gesch\"aftsentwicklung im Berichtsjahr Gegenstand
der Beratungen des Aufsichtsrats waren.\quelleGB{2024}{6} Der Standort
Lauchheim ist seither einer der deutschen Standorte des
Konzerns.\quelleErneut{p11-gb2024-9}

\subsection{Internationalisierung und Kapitalallokation}
\label{subsec:international}

\paragraph{Internationalisierung.}
Der Konzern unterh\"alt \num{\StandorteGesamt} Standorte in den nach eigener
Darstellung drei wichtigsten M\"arkten: die deutschen Standorte Beelen,
Espelkamp, Lauchheim und Limbach-Oberfrohna, die chinesische Gesellschaft in
Changzhou sowie einen Standort in Clayton in den
USA.\quelleErneut{p11-gb2024-9} Die US-Gesellschaft firmiert als Aumann Winding and
Automation Inc., Clayton, und steht zu \Pct{\AnteilUSGesellschaft} im Konzernbesitz.\quelleGB{2024}{52}
Ma{\ss}nahmen zur Reorganisation der Aumann Technologies (China) Ltd.\ waren
Gegenstand der Aufsichtsratsberatungen.\quelleGB{2024}{7}

\paragraph{Kapitalallokation.}
Der Bericht weist zum Jahresende 2024 eine Nettoliquidit\"at von
\MioEUR{\NettoliquiditaetZFIV} und ein Eigenkapital von
\MioEUR{\EigenkapitalZFIV} aus und leitet daraus die M\"oglichkeit ab, sowohl
organisch als auch \"uber gezielte Akquisitionen zu
wachsen.\quelleGB{2024}{5}\quelleMerken{p11-gb2024-5} Vorstand und Aufsichtsrat schlugen eine Dividende
von \EURje{\DividendeJeAktieZFIV} je Aktie vor und legten ein \"offentliches
Aktienr\"uckkaufangebot f\"ur bis zu \num{\RueckkaufAktienMax} eigene Aktien
zu einem Angebotspreis von \EURje{\RueckkaufPreis} je Aktie
auf.\quelleErneut{p11-gb2024-5} Die zur\"uckgekauften Aktien wurden anschlie{\ss}end
eingezogen; die Entwicklung der Aktienzahl ist in
Abschnitt~\ref{subsec:eigentuemer} dargestellt.

\subsection{Bisherige Wirkung auf operatives EBITDA und EBIT}
\label{subsec:wirkung}

Auf Konzernebene stieg das operative EBITDA 2024 um
\Pct{\EbitdaOpDeltaZFIV} auf \MioEUR{\EbitdaOpZFIV}; das EBIT erreichte
\MioEUR{\EbitZFIV}.\quelleGB{2024}{2}\quelleMerken{p11-gb2024-2} Die EBITDA-Marge verbesserte sich von
\Pct{\EbitdaMargeZFIII} auf \Pct{\EbitdaMargeZFIV}.\quelleErneut{p11-gb2024-2}

Der Bericht f\"uhrt die Ergebnisverbesserung im Wesentlichen auf das Segment
E-Mobility zur\"uck, dessen EBITDA sich von \MioEUR{\EbitdaEMobilityZFIII} auf
\MioEUR{\EbitdaEMobilityZFIV} nahezu verdoppelte.\quelleErneut{p11-gb2024-5} Im
Segment Next Automation blieb das EBITDA mit
\MioEUR{\EbitdaNextAutoZFIV} nach \MioEUR{\EbitdaNextAutoZFIII} nahezu
unver\"andert.\quelleErneut{p11-gb2024-13}

Zum Zeitpunkt des Gesch\"aftsberichts 2024 war die Umbenennung des Segments
erst wenige Wochen alt; der Bericht weist der Neuausrichtung daher keinen
eigenen Ergebnisbeitrag zu.

\subsection{Ausblick 2025 und dessen Widerspiegelung in den Berichten}
\label{subsec:ausblick}

Im Gesch\"aftsbericht 2024 prognostizierte die Gesellschaft f\"ur 2025 einen
R\"uckgang des Umsatzes auf \MioEUR{\PrognoseUmsatzZFVUnten} bis
\MioEUR{\PrognoseUmsatzZFVOben} bei einer EBITDA-Marge von
\Pct{\PrognoseMargeZFVUnten} bis \Pct{\PrognoseMargeZFVOben}. Zugleich
erwartete sie erste Anzeichen einer Erholung der Branche im Verlauf des Jahres
2025.\quelleGB{2024}{5}\quelleMerken{p12-gb2024-5}

Diese Prognose wurde im Jahresverlauf unver\"andert best\"atigt: in der
Quartalsmitteilung zum ersten Quartal
2025,\quelleQM{Quartalsmitteilung Q1 2025}{2} im
Halbjahresfinanzbericht\quelleQM{Halbjahresfinanzbericht H1 2025}{3} und in
der Quartalsmitteilung zum dritten Quartal
2025.\quelleQM{Quartalsmitteilung Q3 2025}{3}

\begin{table}[htbp]
\centering
\caption{Prognose f\"ur 2025 aus dem Gesch\"aftsbericht 2024 im Vergleich zum
berichteten Ist. Betr\"age in Mio.\,\euro.}
\label{tab:prognose-ist}
\begin{tabularx}{\linewidth}{@{}Xllc@{}}
\toprule
\textbf{Kennzahl} & \textbf{Prognose GB 2024} & \textbf{Ist 2025} & \textbf{Lage}\\
\midrule
Umsatz
  & \num{\PrognoseUmsatzZFVUnten}--\num{\PrognoseUmsatzZFVOben}
  & \num{\UmsatzZFV}
  & unterhalb\\
EBITDA-Marge (berichtet)
  & \Pct{\PrognoseMargeZFVUnten}--\Pct{\PrognoseMargeZFVOben}
  & \Pct{\EbitdaMargeBerichtetZFV}
  & oberhalb\\
EBITDA-Marge (bereinigt)
  & \Pct{\PrognoseMargeZFVUnten}--\Pct{\PrognoseMargeZFVOben}
  & \Pct{\EbitdaMargeZFV}
  & oberhalb\\
\bottomrule
\end{tabularx}
\end{table}

\noindent Der Umsatz 2025 lag mit \MioEUR{\UmsatzZFV} um
\PctBetrag{\PrognoseAbweichungUmsatz} unter der Untergrenze der
Prognosespanne.\quelleGB{2025}{2}\quelleMerken{p11-gb2025-2} Die berichtete EBITDA-Marge \"ubertraf mit
\Pct{\EbitdaMargeBerichtetZFV} die Obergrenze der Prognosespanne
deutlich.\quelleErneut{p11-gb2025-2} Die Prognose wurde damit beim Umsatz verfehlt und
bei der Marge \"ubertroffen, ohne dass sie im Jahresverlauf angepasst worden
w\"are.
